\section{Finite graphs and the complete hypercube vertex theorem}
\label{sec:hypercube}

This section proves both hypercube inequalities stated in the source article,
including the full vertex minimization theorem.  The arguments operate on
finite bit strings and integer powers.  In particular, the logarithm in the
usual edge statement need not occur in the object language.  The historical
vertex theorem is Harper's \cite{Harper1966,Harper2004}; we give all the
combinatorial steps needed here, together with explicit bounds on the
compression procedure.

\subsection{Finite encodings and graph profiles}

A graph on \(v\) labelled vertices is a symmetric, irreflexive Boolean
\(v\)-by-\(v\) matrix.  A subset is a \(v\)-bit mask.  Its cardinality,
edge boundary, external vertex boundary and closed neighbourhood are,
respectively,
\[
 \begin{aligned}
 &|S|,\qquad N(S)=S\cup\partial_V S,\\
 &\partial_E S=\{\{x,y\}:x\in S,\ y\notin S,\ x\sim y\},\\
 &\partial_V S=\{y\notin S:(\exists x\in S)\ x\sim y\}.
 \end{aligned}
\]
All quantifiers in these definitions are bounded by \(v\), so their
characteristic functions and cardinalities are primitive recursive.
For \(0\leq k\leq v\), the two isoperimetric profiles are computed by
scanning the masks below \(2^v\) of cardinality \(k\) and retaining the
smallest boundary count.  This scan also returns a minimizing mask.
Thus the minima in the article's finite-graph definitions are actual
bounded-search outputs, with no principle of choice required.

Write \(Q_d=\{0,1\}^d\), identifying a vertex with a subset of
\([d]=\{1,\ldots,d\}\).  Adjacent vertices differ in exactly one coordinate.
A vertex is an integer below \(2^d\), and a family \(S\subseteq Q_d\) is
an integer below \(2^{2^d}\).  Put
\[
 k=|S|,\qquad e(S)=\#\{\{x,y\}\subseteq S:x\sim y\},
 \qquad b(S)=|\partial_E S|.
\]
Counting the \(d\) incident edges at each vertex of \(S\) gives the exact
identity
\begin{equation}\label{eq:cube-degree}
 dk=2e(S)+b(S).
\end{equation}
An internal edge is counted twice and an outgoing edge once; this is also a
proof by interchanging two finite sums of Boolean indicators.

\subsection{A logarithm-free sharp edge inequality}

We first record the entirely algebraic form of the elementary
arithmetic--geometric mean inequality used in the proof.

\begin{lemma}[Finite AM--GM]
For nonnegative rational numbers \(x_1,\ldots,x_m\), \(m\geq1\),
\[
 \prod_{j=1}^{m}x_j\ \leq\
 \left(\frac{x_1+\cdots+x_m}{m}\right)^m.
\]
For positive inputs equality holds exactly when all inputs are equal.
\end{lemma}
\begin{proof}
For two inputs, multiply the assertion by \(4\) and use
\((x_1-x_2)^2\geq0\).  Applying this operation to consecutive pairs,
then pairs of pair averages, and so on, proves the assertion for a list
of length \(2^r\).  At every level the average of the list is unchanged;
the product is bounded by the square of the product of the next list.
The product at the root of this finite binary tree is the overall average
raised to \(2^r\).

For an arbitrary positive list of length \(m\), let its average be \(u>0\).
Choose a power \(M=2^r\geq m\), append \(M-m\) copies of \(u\), and
apply the power-of-two result.  The average is still \(u\), so
\((\prod x_j)u^{M-m}\leq u^M\).  Cancelling the positive factor gives
the result.  If an input is zero, the asserted inequality is immediate.
The equality assertion follows from the equality condition at every pair
of the binary tree.  One may choose \(r=m\), so even the padding length
has a specified primitive recursive bound.
\end{proof}

\begin{lemma}[The split inequality]\label{lem:cube-split}
For integers \(a\geq b\geq0\), with the convention \(0^0=1\),
\begin{equation}\label{eq:cube-split}
 4^b a^a b^b\ \leq\ (a+b)^{a+b}.
\end{equation}
If \(a,b>0\), equality holds exactly when \(a=b\).
\end{lemma}
\begin{proof}
The case \(b=0\) is an identity.  For \(b>0\), apply the preceding lemma
to \(b\) entries equal to \(2\) and \(a-b\) entries equal to \(1\):
\[
 2^b\leq\left(\frac{a+b}{a}\right)^a.
\]
Also \(a+b\geq2b\), hence \(2^b\leq((a+b)/b)^b\).
Multiply these inequalities and clear the positive denominators.
If \(a>b>0\), the second inequality is strict; if \(a=b\), both are
equalities.  Thus the equality assertion is decidable integer algebra.
\end{proof}

\begin{theorem}[Hypercube edge inequality, exact arithmetic form]
\label{thm:cube-edge}
For every \(S\subseteq Q_d\),
\begin{equation}\label{eq:cube-edge-internal}
 2^{2e(S)}\leq |S|^{|S|},
\end{equation}
or, equivalently by \eqref{eq:cube-degree},
\begin{equation}\label{eq:cube-edge-integer}
 2^{d|S|}\leq 2^{b(S)}|S|^{|S|}.
\end{equation}
For nonempty \(S\), equality holds exactly when \(S\) is a subcube:
a set obtained by fixing some coordinates and leaving all others free.
\end{theorem}
\begin{proof}
Induct on \(d\).  For \(d=0\) there are only the empty family and the
singleton family, and the assertion is \(1\leq1\).
Split \(S\subseteq Q_d\) into its two sections \(S_0,S_1\subseteq Q_{d-1}\)
in a fixed coordinate.  Let their sizes be \(a,b\), interchanging their
names if necessary so that \(a\geq b\).  Edges between the sections
correspond exactly to the common vertices \(S_0\cap S_1\); put
\(c=|S_0\cap S_1|\leq b\).  Therefore
\[
 e(S)=e(S_0)+e(S_1)+c.
\]
The induction hypotheses and Lemma~\ref{lem:cube-split} give
\[
 2^{2e(S)}
 \leq 2^{2e(S_0)}2^{2e(S_1)}4^b
 \leq a^a b^b4^b
 \leq(a+b)^{a+b}.
\]
This proves the inequality, including either empty section.

For the equality statement, suppose first that both sections are nonempty.
Every inequality in this chain must be an equality.  The split lemma gives
\(a=b\), and equality in \(c\leq b\) gives \(S_0=S_1\).
Both sections have equality in the induction hypothesis, so the common
section is a subcube.  Allowing the split coordinate to be free produces
precisely a subcube in \(Q_d\).  If one section is empty, the nonempty
section has equality by cancellation, so it is a subcube by induction;
the split coordinate is fixed.  Conversely, a subcube with \(r\) free
coordinates has \(k=2^r\), \(2e=r2^r\), and
\(b=(d-r)2^r\).  Substitution proves equality, also for \(r=0\).
\end{proof}

For \(k>0\), the usual real-number reading of
\eqref{eq:cube-edge-internal} is
\[
 2e(S)\leq k\log_2 k,
 \qquad b(S)\geq k(d-\log_2 k).
\]
The exact finite theorem proved inside arithmetic is
\eqref{eq:cube-edge-integer}; this avoids using a real logarithm as an
undefined arithmetic object.  The equality witnesses in the article are
therefore obtained constructively, and the argument in fact classifies
all nonempty equality cases.

\begin{corollary}[Sharp expansion at half volume]
For \(d\geq1\) and \(0<|S|\leq2^{d-1}\), one has
\(b(S)\geq |S|\).  This constant is sharp, as witnessed by a coordinate
half-cube.  Applying the same assertion to the complement gives
\[
 b(S)\geq\min\{|S|,2^d-|S|\}.
\]
\end{corollary}
\begin{proof}
If \(b(S)<k\), then \eqref{eq:cube-edge-integer} and \(k\leq2^{d-1}\)
would imply
\[
 2^{dk}\leq2^{b(S)}k^k
 \leq2^{b(S)+(d-1)k}<2^{dk},
\]
an impossible comparison of integers.  Since integer order is decidable,
this proves \(b(S)\geq k\) intuitionistically.  A coordinate half-cube
has exactly one outgoing edge at every vertex.  Complementation preserves
the edge boundary.
\end{proof}

\subsection{Simplicial order and its neighbourhoods}

Order subsets of \([d]\) first by cardinality.  Within one cardinality,
use lexicographic order with a smaller coordinate preferred: if
\(X\neq Y\) have equal size, then \(X<Y\) when the least element of
\(X\mathbin{\triangle}Y\) belongs to \(X\).
For example, the two-element subsets of \([4]\) are ordered
\[
 12,\ 13,\ 14,\ 23,\ 24,\ 34.
\]
This is the \emph{simplicial order}.  Let \(I_d(k)\) denote its first
\(k\) vertices, and let \(\operatorname{rk}_d(X)\) be the rank of
\(X\), starting at zero.  Both are primitive recursive: enumerate the
\(2^d\) masks, count their bits, compare them in the stated order, and
count how many precede the given mask.

\begin{lemma}[Neighbourhoods preserve initial segments]
\label{lem:cube-initial-neighborhood}
The closed neighbourhood of \(I_d(k)\) is another simplicial initial
segment.  Moreover, \(k\mapsto |N(I_d(k))|\) is nondecreasing.
\end{lemma}
\begin{proof}
The empty and full initial segments are immediate.  Any other nonempty
initial segment consists of all sets of size less than \(r\), together
with a nonempty lexicographic initial family \(F\) of \(r\)-sets.
Its neighbourhood contains every set of size at most \(r\): for a set
of size \(r\) this follows by deleting one element, unless \(r=0\),
which is immediate.  Its only other members are the \((r+1)\)-sets
containing some member of \(F\).

If \(Y=\{y_1<\cdots<y_{r+1}\}\), its lexicographically first
\(r\)-subset is \(\{y_1,\ldots,y_r\}\).  Consequently \(Y\) contains
a member of \(F\) exactly when this first \(r\)-subset belongs to \(F\).
Deleting the largest element preserves the weak lexicographic order on
sets of the same size: compare their sorted lists at the first differing
position; if this is the last position, the shortened lists agree.
It follows that the eligible \((r+1)\)-sets form an initial segment.
For \(r=0\), every singleton contains the empty set, giving the same
conclusion.  Finally \(I_d(k)\subseteq I_d(k+1)\), and the operation
\(N\) preserves inclusion.
\end{proof}

\subsection{Compression with a specified termination bound}

For coordinate \(i\), write \(S_i^0,S_i^1\subseteq Q_{d-1}\) for the
sections obtained by fixing that coordinate to \(0\) or \(1\) and
then deleting it.  The \(i\)-compression \(C_iS\) replaces these two
sections by
\[
 I_{d-1}(|S_i^0|),\qquad I_{d-1}(|S_i^1|).
\]
Deleting a coordinate from vertices which agree in that coordinate
preserves their relative simplicial order.  Thus each section is moved
leftwards in its inherited order, without changing its size.

\begin{lemma}[One compression step]
\label{lem:cube-compression}
Assume the vertex theorem in dimension \(d-1\): every family \(T\)
in that cube satisfies
\( |N(T)|\geq|N(I_{d-1}(|T|))|\).
Then \( |C_iS|=|S|\) and
\[
 |N(C_iS)|\leq |N(S)|.
\]
\end{lemma}
\begin{proof}
Put \(A=S_i^0\), \(B=S_i^1\), and let \(A',B'\) be their compressed
initial segments.  A vertex in the zero section of \(N(S)\) either
is in the closed neighbourhood of \(A\) inside that section or is
joined across coordinate \(i\) to a member of \(B\).  Hence the two
sections of \(N(S)\) are
\[
 N(A)\cup B,\qquad N(B)\cup A.
\]
The corresponding compressed sections are \(N(A')\cup B'\) and
\(N(B')\cup A'\).  By
Lemma~\ref{lem:cube-initial-neighborhood}, \(N(A')\) and \(B'\)
are nested initial segments.  Therefore
\[
 \begin{split}
 |N(A')\cup B'|
 &=\max\{|N(A')|,|B'|\}\\
 &\leq\max\{|N(A)|,|B|\}
 \leq|N(A)\cup B|.
 \end{split}
\]
The first inequality uses the dimension induction hypothesis and
\(|B'|=|B|\).  Interchange \(A,B\) for the other section and add.
The cardinality assertion is part of the definition of compression.
\end{proof}

Define the nonnegative integer potential
\[
 P(S)=\sum_{X\in S}\operatorname{rk}_d(X).
\]
If \(C_iS\neq S\), at least one section has moved strictly leftwards;
matching its sorted old and new lists shows that its rank sum strictly
decreases.  The other section's rank sum does not increase.  Thus
\[
 P(C_iS)<P(S)\qquad(C_iS\neq S),
 \qquad P(S)\leq |S|(2^d-1)<2^{2d}.
\]
For a total algorithm, choose the least coordinate whose compression
changes the current family, and leave the family unchanged if there is
none.  Iterate this instruction \emph{exactly} \(2^{2d}\) times.
Each nontrivial step lowers the potential by at least one, and once the
family is fixed, all later steps leave it fixed.  The final family \(B\)
is therefore fixed by every coordinate compression.  This is a bounded
primitive recursive iteration, not an unbounded ``repeat until stable''
search.  Lemma~\ref{lem:cube-compression} gives
\[
 |B|=|S|,\qquad |N(B)|\leq|N(S)|.
\]

\subsection{The fully compressed exceptional families}

Being fixed by every coordinate compression does not by itself make a
family an initial segment.  We now handle exactly the possible obstruction;
omitting it would leave the proof incomplete.

\begin{lemma}[Classification of a compressed obstruction]
\label{lem:cube-exceptions}
Suppose \(B\subseteq Q_d\) is fixed by every coordinate compression and
is not an initial segment.  Let \(x\) be its first missing vertex and
\(y\) its last present vertex.  Then \(x<y\), \(x,y\) are consecutive
in simplicial order, and \(y=[d]\setminus x\).  In particular, \(B\)
is the initial segment of size \(|B|\) with \(x\) replaced by \(y\).
There are only the following two forms:
\begin{enumerate}
\item If \(d=2r+1\), then the initial segment is all sets of size at most
\(r\), and
\[
 x=\{r+2,\ldots,2r+1\},\qquad y=\{1,\ldots,r+1\}.
\]
\item If \(d=2r\), then the initial segment consists of all sets of size
less than \(r\), followed by all \(r\)-sets containing \(1\), and
\[
 x=\{1,r+2,\ldots,2r\},\qquad y=\{2,\ldots,r+1\}.
\]
\end{enumerate}
\end{lemma}
\begin{proof}
Take any inversion \(u<v\) with \(u\notin B\) and \(v\in B\).
If they agreed in some coordinate, they would belong to the same section;
that section is an initial segment, contradicting the inversion.
They must therefore disagree in every coordinate, so \(v=[d]\setminus u\).
Apply this to \(x,y\).  If a vertex \(z\) lay strictly between them,
then \(z\in B\) would force \(z=[d]\setminus x=y\), whereas
\(z\notin B\) would force \(z=[d]\setminus y=x\).
Both alternatives are impossible.  Membership is decidable, so this is
an intuitionistically legitimate case distinction.  Thus \(x,y\) are
consecutive, and the definition of first missing and last present gives
the claimed single exchange.

Put \(p=|x|\) and \(q=|y|\).  Since \(x<y\) and they are
complements, \(p\leq q\) and \(p+q=d\).
If \(p<q\), consecutiveness forces \(q=p+1\); otherwise an entire
intermediate cardinality lies between them.  The vertex \(x\) must be
the last \(p\)-set, and \(y\) the first \((p+1)\)-set.  These are the
displayed odd-dimensional vertices.
If \(p=q=r\), then \(d=2r\).  Because they are complementary, their
least differing coordinate is \(1\), which belongs to \(x\).
All \(r\)-sets containing \(1\) precede every \(r\)-set not containing
\(1\).  Consecutiveness therefore forces \(x\) to be the last of the
former kind and \(y\) the first of the latter kind, giving the second
form.
\end{proof}

\begin{lemma}[Repairing the obstruction]
\label{lem:cube-repair}
For a family \(B\) in Lemma~\ref{lem:cube-exceptions}, put
\(I=I_d(|B|)\).  Then \(|N(I)|\leq |N(B)|\).
\end{lemma}
\begin{proof}
In the odd case \(d=2r+1\), suppose first \(r\geq1\).
The neighbourhood \(N(I)\) consists of all sets of size at most \(r+1\).
The family \(B\) still contains every set of size at most \(r-1\),
so every set of size at most \(r\) belongs to \(N(B)\).
Every \((r+1)\)-set has \(r+1\geq2\) different \(r\)-subsets.
Only one \(r\)-set was removed; hence at least one of these subsets
remains in \(B\).  Consequently \(N(I)\subseteq N(B)\).
For \(r=0\), both singleton families in \(Q_1\) have the entire
\(Q_1\) as their closed neighbourhood.

In the even case \(d=2r\), suppose \(r\geq2\).
Here \(N(I)\) contains all sets of size at most \(r\), and all
\((r+1)\)-sets containing \(1\).  The smaller sets still belong to
\(N(B)\) because all sets of size less than \(r\) remain in \(B\).
Each \((r+1)\)-set containing \(1\) has \(r\geq2\) different
\(r\)-subsets containing \(1\), and at most one was removed.
Again \(N(I)\subseteq N(B)\).
For \(r=1\), the two families are \(\{\varnothing,\{1\}\}\)
and \(\{\varnothing,\{2\}\}\) in \(Q_2\); both neighbourhoods
are all of \(Q_2\).
\end{proof}

\begin{theorem}[Harper's vertex theorem]\label{thm:cube-harper}
For every \(d\geq0\) and every \(S\subseteq Q_d\),
\begin{equation}\label{eq:cube-harper}
 |N(S)|\geq |N(I_d(|S|))|.
\end{equation}
Consequently, for every \(0\leq k\leq2^d\),
\[
 \Phi_V(Q_d,k)=|N(I_d(k))|-k.
\]
\end{theorem}
\begin{proof}
Induct on \(d\), starting with the two families in \(Q_0\).
The induction hypothesis justifies every compression step by
Lemma~\ref{lem:cube-compression}.  Run the bounded iteration to obtain
\(B\) fixed by all compressions.  If \(B\) is an initial segment,
then \(B=I_d(|S|)\).  Otherwise
Lemmas~\ref{lem:cube-exceptions} and \ref{lem:cube-repair} give
\[
 |N(I_d(|S|))|\leq|N(B)|\leq|N(S)|.
\]
Whether a finite family is an initial segment is decidable by a bounded
scan, so the displayed case split is constructive.
Subtract the common cardinality \(k\) from the closed-neighbourhood
inequality to compare external boundaries; the initial segment itself
is a minimizing witness.
\end{proof}

\begin{corollary}[The binomial thresholds in the source article]
\label{cor:cube-thresholds}
Let \(0\leq r\leq d\), and put
\[
 B(d,r)=\sum_{j=0}^{r}\binom dj,
 \qquad \binom d{d+1}=0.
\]
If \(|S|\geq B(d,r)\), then
\[
 |N(S)|\geq B(d,\min\{r+1,d\}).
\]
At the exact threshold,
\[
 \Phi_V(Q_d,B(d,r))=\binom d{r+1}.
\]
\end{corollary}
\begin{proof}
The initial segment of size \(B(d,r)\) is the exact Hamming ball
\(\{X:|X|\leq r\}\).  Its closed neighbourhood is the ball of radius
\(\min\{r+1,d\}\).  This follows by adding or deleting one element;
its cardinality is the stated binomial sum.  Harper's theorem and
monotonicity of the neighbourhoods of initial segments give the first
claim.  At equality of cardinalities, subtract \(B(d,r)\) and use the
ball as a witness.
\end{proof}

For intermediate sizes, the minimizing family consists of all layers below
some weight and an \emph{initial lexicographic portion} of the next layer.
The qualification on the partial layer matters: the proof does not assert
that every arbitrarily chosen portion of that layer is optimal.
The same argument also yields all-distance neighbourhoods.  If \(N^t\)
denotes \(t\) iterations of \(N\), then
\[
 |N^t(S)|\geq |N^t(I_d(|S|))|.
\]
Indeed, apply Harper at each step, use monotonicity in the cardinality,
and use Lemma~\ref{lem:cube-initial-neighborhood} to identify each next
initial segment.  At a binomial threshold this gives the ball of radius
\(\min\{r+t,d\}\).  The induction is a bounded iteration if \(t\)
is replaced by \(\min\{t,d\}\), since every nonempty neighbourhood
has filled the cube by distance \(d\).

\subsection{Why these proofs use only intuitionistic primitive recursion}

Here is an explicit audit of the induction strength.  For the dimension
induction in Harper's theorem, use the Boolean predicate
\[
 H(d)=\bigwedge_{m<2^{2^d}}
 \bigl[\,|N(S_m)|\geq|N(I_d(|S_m|))|\,\bigr],
\]
where \(S_m\) is the family encoded by mask \(m\), and the conjunction
is evaluated by a primitive recursive finite loop.  The induction formula
is simply \(H(d)=1\), a quantifier-free equation.  The same construction
works for the edge inequality and for its equality classification.
Intermediate universal assertions range over coordinates, vertices,
families, or time indices with the displayed primitive recursive bounds;
replace each bounded quantifier by its finite Boolean fold.

The potential argument uses at most \(2^{2d}\) states of a deterministic
loop.  Each state is a mask below \(2^{2^d}\), each rank is below \(2^d\),
and each potential is below \(2^{2d}\).  A code for the whole computation
trace can, for example, be bounded by
\[
 (2^{2^d})^{\,2^{2d}+1}.
\]
This generous bound suffices to compile every invariant to a bounded
arithmetical predicate.  Exact counting, comparison, addition,
multiplication, exponentiation, bit extraction, binomial coefficients and
finite sorting are all primitive recursive.

For the split inequality, the averaging proof only uses finitely many
rational operations.  In its application the inputs are \(1,2\), or their
average \((a+b)/a\); all averages lie between \(1\) and \(2\).
With padding length \(M=2^a\), a common denominator \(aM\) represents
all intermediate means, with numerators at most \(2aM\).  Thus the
binary averaging tree also has an explicit primitive recursive code
bound.  Clearing positive denominators leaves only integer polynomial
inequalities.  No unbounded search, unbounded induction formula,
nonconstructive extremum, or real-number case distinction is used.

The supplied \texttt{hypercube\_verify.py} implements the exact masks,
boundaries, simplicial initial segments and bounded compression algorithm.
Its recorded run checks both main inequalities on every family in
\(Q_d\) for \(0\leq d\leq4\), checks the binomial thresholds, and
checks that the nonempty equality cases of the edge inequality number
\(3^d\), exactly the number of subcubes.  These finite checks are useful
regression tests for the definitions and the exceptional cases.  They are
not the uniform proof: that proof is the quantified, bounded-induction
argument given above.
